\chapter{Background and Related Work}\label{ch:background}

\section{TCP}

TCP~\cite{rfc9293} offers a reliable, ordered byte stream between two endpoints. Reliability is built from four
mechanisms:
\begin{itemize}
  \item cumulative acknowledgements, optionally extended with selective acknowledgements (SACK)~\cite{rfc2018};
  \item a retransmission timer computed from a smoothed round-trip time and its variance~\cite{rfc6298};
  \item fast retransmit on duplicate acknowledgements;
  \item window scaling for large bandwidth-delay products~\cite{rfc7323}.
\end{itemize}
Congestion control follows the additive-increase, multiplicative-decrease principle~\cite{jacobson,chiu-jain}:
slow start doubles the window every round trip until the first loss, after which the window grows by about one
segment per round trip and halves on loss~\cite{rfc5681}. The steady-state throughput of such a flow is roughly
proportional to $\mathit{MSS}/(\mathit{RTT}\sqrt{p})$ for loss rate $p$~\cite{mathis}, which is why loss and
latency together dominate bulk performance on long paths.

Two properties make TCP a poor fit for mixed traffic. First, the byte stream is totally ordered: when a segment
is lost, the bytes after it cannot be delivered to the application until the retransmission arrives, even if
they belong to an unrelated message. Second, the whole connection shares one congestion window, so the
behaviour of the bulk data determines the latency of everything else.

\section{QUIC}

QUIC~\cite{rfc9000,quic-design} runs over UDP and multiplexes many independent streams inside one encrypted
connection. A lost packet only blocks the streams whose data it carried. QUIC has two further features worth
noting:
\begin{itemize}
  \item \emph{packet numbers that are never reused}, so retransmitted data travels in a new packet and RTT
    samples are unambiguous;
  \item \emph{address validation}: a Retry token proves that a client can receive at its claimed address before
    the server commits resources or sends large responses.
\end{itemize}
Loss detection and congestion control are specified in RFC~9002~\cite{rfc9002}, and packet protection with
TLS~1.3 in RFC~9001~\cite{rfc9001,rfc8446}. RFC~9001 also bounds how many packets a single AEAD key may protect.

For the problem in this report, QUIC still has three drawbacks:
\begin{itemize}
  \item Congestion control is per connection, so a bulk stream's losses throttle the latency-sensitive stream
    that shares the connection.
  \item Encryption is mandatory, so a trusted link always pays for TLS.
  \item Reliability is a property of the connection's streams. Unreliable delivery requires the separate
    datagram extension and is not a per-stream setting.
\end{itemize}
Several of Fuse's version~8 mechanisms deliberately borrow QUIC's well-studied answers: Retry-style address
tokens, path validation, AEAD key limits and the separation of RTT estimation from retransmission.

\section{Other transports}

\begin{itemize}
  \item \textbf{SCTP}~\cite{sctp} provides multiple streams within an association, with per-stream ordering and
    an optional unordered delivery mode. It is message-oriented, like Fuse's SDK. Its congestion control,
    however, is per association, and in practice it is often blocked by middleboxes.
  \item \textbf{UDT}~\cite{udt} is a UDP-based protocol for high-speed bulk transfer over wide-area networks,
    with its own rate-based congestion control. It addresses a single bulk flow rather than mixed traffic.
  \item \textbf{DTLS}~\cite{rfc6347,rfc9147} adapts TLS to datagrams. Its record layer keeps per-session
    sequence numbers, an epoch and a replay window. As Section~\ref{sec:crypto} explains, that state is why
    Fuse does not use one DTLS session per lane for its high-throughput path.
  \item \textbf{TFTP}~\cite{rfc1350} has the server answer each request from a fresh port, so that one
    well-known port can serve many transfers without connection identifiers. Fuse's SDK reuses this idea.
\end{itemize}

\section{Systems techniques relied upon}

\begin{itemize}
  \item \textbf{Batched UDP I/O.} One system call per datagram limits a modern host to a few hundred thousand
    datagrams per second, however fast the network. Linux provides \code{recvmmsg}/\code{sendmmsg} to move many
    datagrams per call~\cite{recvmmsg}, and generic segmentation offload (\code{UDP\_SEGMENT}, Linux 4.18): the
    application hands the kernel one large buffer and the kernel or network card splits it into ordinary
    MTU-sized datagrams~\cite{gso}.
  \item \textbf{Path MTU.} Datagrams larger than the path MTU are fragmented at the IP layer, and the loss of any
    fragment destroys the whole datagram. Packetization-layer path MTU discovery~\cite{rfc8899} probes larger
    sizes carefully and backs off after failure. Fuse's adaptive block size follows its spirit.
  \item \textbf{Authenticated encryption.} AES-GCM~\cite{nist-gcm} provides confidentiality and integrity in one
    pass. Its security collapses if a (key, nonce) pair is ever reused with different inputs: two such
    ciphertexts allow the authentication key to be recovered and arbitrary messages to be
    forged~\cite{joux}. Every encryption design in Fuse is shaped by this constraint.
  \item \textbf{Key derivation and keyed hashing.} HKDF~\cite{rfc5869} derives keys from a pre-shared secret and
    a salt. SipHash~\cite{siphash} is a fast keyed pseudorandom function for short inputs, suitable for
    unforgeable tokens.
  \item \textbf{Lock-free hand-off.} A ring buffer with one producer and one consumer needs no locks, only
    correctly ordered reads and writes of its two indices~\cite{lamport-spsc}.
\end{itemize}

\section{Summary of the gap}

Table~\ref{tab:compare} summarises the comparison that motivates Fuse.

\begin{table}[H]
\centering\small
\caption{Transport properties relevant to mixed loss-tolerant and loss-intolerant traffic.}\label{tab:compare}
\begin{tabularx}{\linewidth}{>{\raggedright\arraybackslash}p{3.6cm}XXX}
\toprule
 & \textbf{TCP} & \textbf{QUIC} & \textbf{Fuse} \\
\midrule
Head-of-line blocking & across all data & per stream & per stream \\
Congestion control & per connection & per connection & per stream \\
Reliability & always & always (datagrams via extension) & per stream (\code{LOSSLESS}) \\
Ordering & always & per stream & per stream (\code{ORDERED}) \\
Encryption & optional (TLS) & mandatory (TLS~1.3) & optional, local (PSK) \\
Forward secrecy & with TLS & yes & DTLS path only (not on the PSK path) \\
Batched I/O & kernel & implementation-dependent & GSO + \code{recvmmsg} \\
\bottomrule
\end{tabularx}
\end{table}
